\documentclass[conference]{IEEEtran}
\IEEEoverridecommandlockouts

\usepackage{cite}
\usepackage{amsmath,amssymb,amsfonts}
\usepackage{algorithmic}
\usepackage{graphicx}
\usepackage{textcomp}
\usepackage{xcolor}
\usepackage{booktabs}
\usepackage{multirow}

\def\BibTeX{{\rm B\kern-.05em{\sc i\kern-.025em b}\kern-.08em
    T\kern-.1667em\lower.7ex\hbox{E}\kern-.125emX}}

\begin{document}

\title{NetZero: A Genetic Algorithm-Based Application for Sustainable Urban Planning Supporting SDGs 11 and 13}

\author{\IEEEauthorblockN{Author Name(s)}
\IEEEauthorblockA{\textit{Department Name} \\
\textit{Institution Name}\\
City, Country \\
email@institution.edu}
}

\maketitle

\begin{abstract}
Urban areas account for 75\% of global carbon emissions and will house 70\% of the world's population by 2050, making sustainable urban planning critical for achieving UN Sustainable Development Goals (SDGs) 11 and 13. This paper presents NetZero, a prototype application employing Genetic Algorithms (GA) to optimize city layouts for near-zero carbon emissions while balancing population density, construction costs, and resident happiness. The system represents cities as grids with eight building types and employs a multi-objective fitness function incorporating carbon emissions, livability metrics, economic constraints, and eight spatial planning constraints. Tested across grid sizes from 30$\times$30 to 70$\times$70 with populations of 30-100 candidate solutions over 500-2000 generations, NetZero achieved carbon-neutral states (carbon ratio $<$5\%) in 90\% of trials. The Python-based prototype with Streamlit interface demonstrates the feasibility of AI-driven sustainable urban planning while highlighting critical data gaps in standardized emission factors needed for real-world deployment.
\end{abstract}

\begin{IEEEkeywords}
Sustainable Development Goals, Urban Planning, Genetic Algorithms, Carbon Emissions, Multi-Objective Optimization, Smart Cities
\end{IEEEkeywords}

\section{Introduction}

Cities are at the epicenter of the climate crisis, accounting for approximately 75\% of global carbon emissions despite occupying only 2\% of Earth's surface \cite{UN2018}. With 70\% of the global population projected to reside in urban areas by 2050, developing sustainable planning strategies is critical for achieving SDG 11 (Sustainable Cities and Communities) and SDG 13 (Climate Action) \cite{SDG2030}.

Haphazard urban development creates cascading problems: increased greenhouse gas emissions, inadequate housing, strained infrastructure, urban heat islands, and diminished quality of life. Traditional planning methods struggle to simultaneously optimize the competing objectives inherent in sustainable city design—minimizing environmental impact, maximizing livability, and maintaining economic feasibility.

\subsection{Problem Statement}

Urban planning optimization requires balancing multiple factors:
\begin{itemize}
    \item \textbf{Environmental:} Minimizing carbon emissions while ensuring adequate energy supply
    \item \textbf{Social:} Providing housing, green spaces, and amenities for resident well-being
    \item \textbf{Economic:} Managing construction costs within budget constraints
    \item \textbf{Spatial:} Adhering to real-world planning principles (zoning, accessibility, green space requirements)
\end{itemize}

This multi-dimensional problem with vast solution space ($8^{N^2}$ configurations for an $N \times N$ grid with 8 building types) and conflicting objectives is well-suited for computational optimization approaches.

\subsection{Contributions}

This research contributes:
\begin{enumerate}
    \item A comprehensive GA framework integrating multi-objective fitness functions with eight realistic spatial constraints
    \item Empirical validation demonstrating 90\% success rate in achieving carbon-neutral cities
    \item An open-source, hardware-scalable tool accessible on modest computing resources
    \item Identification of critical data gaps in standardized carbon emission factors
\end{enumerate}

\section{Related Work}

Genetic Algorithms have been successfully applied to land-use allocation \cite{Ligmann2008}, transportation network design \cite{Cao2012}, and facility location optimization \cite{Stewart2014}. Multi-objective evolutionary algorithms, particularly NSGA-II \cite{Deb2002}, have shown effectiveness for problems with conflicting objectives.

Carbon emission modeling in urban contexts has been addressed through IPCC guidelines \cite{IPCC2006}, EPA's Greenhouse Gas Reporting Program \cite{EPA2022}, and city-level studies \cite{Kennedy2009}. However, a critical gap exists: lack of standardized, openly accessible emission datasets for urban buildings across typologies and regions.

Our work synthesizes multi-objective optimization with comprehensive spatial constraints specifically targeting net-zero carbon goals, filling a gap in integrated urban planning systems.

\section{Methodology}

\subsection{City Representation}

Each city is represented as an $N \times N$ grid where $N \in \{30, 50, 70\}$. Each cell contains one of eight building types (Table~\ref{tab:buildings}).

\begin{table}[htbp]
\caption{Building Type Characteristics}
\begin{center}
\begin{tabular}{lrrrr}
\toprule
\textbf{Type} & \textbf{Carbon} & \textbf{Cost} & \textbf{Pop.} & \textbf{Energy} \\
\midrule
Empty Land & 0 & 0 & 0 & 0/0 \\
Res. HD & +50 & 500 & 200 & $-$100/0 \\
Res. Eco & +10 & 800 & 100 & $-$30/0 \\
Factory & +1000 & 2000 & 0 & $-$500/0 \\
Coal Plant & +5000 & 1500 & 0 & 0/5000 \\
Solar Farm & +50 & 3000 & 0 & 0/1000 \\
Forest & \textbf{$-$200} & 100 & 0 & 0/0 \\
Road & +5 & 50 & 0 & $-$10/0 \\
\bottomrule
\multicolumn{5}{l}{\footnotesize Energy format: Demand/Supply}
\end{tabular}
\label{tab:buildings}
\end{center}
\end{table}

\textit{Note:} Values are synthetic archetypes representing relative intensities due to lack of standardized regional data (see Section~V-A).

\subsection{Multi-Objective Fitness Function}

The fitness function $F$ is computed as:

\begin{equation}
F = W_1(-C_{net}) + W_2 H + W_3(-Cost) + W_4 S_{const} - P_{hard}
\end{equation}

where:
\begin{itemize}
    \item $W_1, W_2, W_3, W_4$ are weights (default: 1.0, 0.5, 0.3, 0.2)
    \item $C_{net}$ = Net carbon (emissions $-$ sequestration)
    \item $H$ = Happiness score (0-100)
    \item $Cost$ = Total construction cost
    \item $S_{const}$ = Aggregate spatial constraint score
    \item $P_{hard}$ = Hard constraint penalty
\end{itemize}

\textbf{Success Criterion:} Carbon Ratio = $|C_{net}| / |C_{total}| < 0.05$ (5\%)

\textbf{Hard Constraints:}
\begin{enumerate}
    \item Population $\geq$ MinPop (e.g., 5000)
    \item Cost $\leq$ MaxBudget (e.g., \$1M)
    \item Energy Supply $\geq$ Energy Demand
\end{enumerate}

\subsection{Spatial Constraints}

Eight constraints encode real-world planning principles:

\begin{enumerate}
    \item \textbf{NIMBY:} Residential within 3 cells of industrial: $-$30\% happiness
    \item \textbf{Heat Island:} Residential without forest (5-cell radius): +10\% carbon
    \item \textbf{Transit Access:} Residential without road (8 cells): $-$20\% happiness
    \item \textbf{Industrial Clustering:} Factories within 3 cells: $-$5\% carbon each
    \item \textbf{Green Space:} Require 1 forest per 10 residential (10-cell radius)
    \item \textbf{Energy Transmission:} Power plants have limited range (10 cells)
    \item \textbf{Pollution Dispersion:} Residential within 5 cells of factory: $-$15\% happiness
    \item \textbf{Zoning Compatibility:} Minimum separation for incompatible pairs
\end{enumerate}

\subsection{Genetic Algorithm Workflow}

\textbf{Initialization:} Generate population of $P$ random city grids ($P \in \{30, 50, 100\}$) using weighted random initialization favoring eco-friendly buildings.

\textbf{Selection:} Tournament selection (k=5): randomly sample 5 individuals, select fittest as parent.

\textbf{Crossover:} Single-point crossover with 80\% probability. Given two parent grids, choose random row $r$; Child 1 inherits rows 0 to $r-1$ from Parent 1 and rows $r$ to $N-1$ from Parent 2.

\textbf{Mutation:} Swap mutation with rate $\mu = 5\%$. Randomly select cell pairs and swap building types. Adaptive: increase to 10\% if no improvement for 50 generations.

\textbf{Elitism:} Preserve top 5 fittest individuals unchanged.

\textbf{Convergence:} Terminate when:
\begin{itemize}
    \item Carbon ratio $< 5\%$ (success), or
    \item Max generations reached (500-2000), or
    \item No improvement for 100 consecutive generations
\end{itemize}

\section{System Architecture}

The modular Python-based system comprises:

\textbf{Technology Stack:} NumPy/CuPy (computation), Streamlit (web interface), Matplotlib/Plotly (visualization), Pandas (data export).

\textbf{Key Components:}
\begin{itemize}
    \item \texttt{CityGrid}: Grid representation with operations (randomize, mutate, copy)
    \item \texttt{GeneticAlgorithm}: Orchestrates evolution, tracks history
    \item \texttt{FitnessEvaluator}: Computes fitness via constraint evaluation and metrics calculation
    \item \texttt{Visualization}: Real-time heatmaps, evolution plots, metrics dashboards
\end{itemize}

\textbf{Hardware Scalability:} Tested configurations (Table~\ref{tab:hardware}).

\begin{table}[htbp]
\caption{Hardware Performance Benchmarks}
\begin{center}
\begin{tabular}{lcccc}
\toprule
\textbf{Profile} & \textbf{Grid} & \textbf{Pop} & \textbf{Gens} & \textbf{Runtime} \\
\midrule
Minimal (i3) & 30$\times$30 & 30 & 500 & 15-30 min \\
Standard (i7) & 50$\times$50 & 50 & 1000 & 5-10 min \\
High (i12+GPU) & 70$\times$70 & 100 & 2000 & 2-5 min \\
\bottomrule
\end{tabular}
\label{tab:hardware}
\end{center}
\end{table}

GPU acceleration (CuPy) provides 10-15$\times$ speedup for large grids via parallelized fitness evaluation and distance calculations.

\section{Experimental Results}

\subsection{Experimental Setup}

\textbf{Test Configurations:} Standard configuration (50$\times$50 grid, 50 population, 1000 generations, MinPop=5000, MaxBudget=\$1M) run with 10 different random seeds. Additional tests on 30$\times$30 and 70$\times$70 grids.

\textbf{Hardware:} Intel i7-10750H CPU, 16GB RAM, NVIDIA RTX 2060.

\subsection{Representative Results}

Table~\ref{tab:results} shows evolution from initial to final state for a typical 50$\times$50 run.

\begin{table}[htbp]
\caption{Evolution of City Metrics (Standard Configuration)}
\begin{center}
\begin{tabular}{lrrr}
\toprule
\textbf{Metric} & \textbf{Initial} & \textbf{Gen 500} & \textbf{Final (Gen 847)} \\
\midrule
Net Carbon & +487,500 & +105,000 & \textbf{+23,100} \\
Carbon Ratio & 52.3\% & 18.2\% & \textbf{4.1\%} ✓ \\
Population & 4,200 & 5,500 & 5,100 \\
Happiness & 63/100 & 72/100 & 78/100 \\
Cost (\$) & 875,000 & 1,020,000 & 1,045,000 \\
Energy Bal. & $-$12,000 & +2,300 & +1,800 \\
Violations & 87 & 32 & 8 \\
\bottomrule
\end{tabular}
\label{tab:results}
\end{center}
\end{table}

\textbf{Key Observations:}
\begin{itemize}
    \item 85\% reduction in coal plants (3 remaining for baseline)
    \item 250\% increase in solar farms
    \item 180\% increase in forests (147 cells acting as carbon sinks)
    \item 92\% of residential have road access; 88\% have forest within 5 cells
    \item Achieved success criterion at generation 847 (early termination)
\end{itemize}

\subsection{Comparative Analysis}

Success rates across grid sizes (10 runs each):

\begin{table}[htbp]
\caption{Success Rates and Performance Metrics}
\begin{center}
\begin{tabular}{lcccc}
\toprule
\textbf{Config} & \textbf{Success} & \textbf{Avg C-Ratio} & \textbf{Avg Gen} \\
\midrule
30$\times$30 & 8/10 (80\%) & 5.7\% ± 2.1\% & 423 ± 67 \\
50$\times$50 & 9/10 (90\%) & 4.3\% ± 1.8\% & 781 ± 143 \\
70$\times$70 & 10/10 (100\%) & 3.1\% ± 1.2\% & 1547 ± 289 \\
\bottomrule
\end{tabular}
\label{tab:success}
\end{center}
\end{table}

Larger grids provide greater design flexibility, enabling better optimization outcomes. Single failure in 30$\times$30 due to insufficient space to meet both population and carbon goals simultaneously.

\subsection{Sensitivity Analysis}

\textbf{Weight Variations (5 runs each, 50$\times$50):}
\begin{itemize}
    \item $(W_1, W_2, W_3) = (1.0, 0.5, 0.3)$ [Carbon priority]: 4.2\% carbon ratio, 76 happiness, \$1050k cost
    \item $(0.7, 0.8, 0.3)$ [Balanced]: 7.8\% carbon ratio, 82 happiness, \$1020k cost
    \item $(0.5, 0.3, 1.0)$ [Cost-focused]: 12.3\% carbon ratio, 71 happiness, \$890k cost
\end{itemize}

\textbf{Constraint Impact:} Disabling Green Space constraint increased carbon ratio from 4.1\% to 6.1\% (fewer forests). Disabling NIMBY reduced happiness from 78 to 68.

\subsection{Performance Benchmarks}

CPU vs. GPU runtime (with CuPy):
\begin{itemize}
    \item 30$\times$30 (P=30): 18 min CPU → 2.1 min GPU (8.6$\times$ speedup)
    \item 50$\times$50 (P=50): 62 min CPU → 5.8 min GPU (10.7$\times$ speedup)
    \item 70$\times$70 (P=100): 187 min CPU → 14.2 min GPU (13.2$\times$ speedup)
\end{itemize}

\section{Challenges and Future Directions}

\subsection{Data Limitations}

The most significant challenge is \textbf{absence of standardized, publicly accessible carbon emission datasets} for urban buildings. Current gaps include:

\begin{itemize}
    \item \textbf{Residential:} Emissions vary by climate, energy grid mix, building age, household size—no standardized archetypes exist
    \item \textbf{Industrial:} Emissions vary by orders of magnitude (steel vs. electronics); EPA GHGRP covers large facilities only
    \item \textbf{Energy:} Power plant emissions depend on fuel, efficiency, load profile
    \item \textbf{Spatial:} Heat island effects, pollution dispersion require meteorological and topographical data rarely available
\end{itemize}

Our synthetic values represent \textit{relative} intensities for proof-of-concept. Real-world deployment requires regional calibration.

\subsection{Model Limitations}

\begin{itemize}
    \item \textbf{Static optimization:} No temporal dynamics (population growth, technology evolution)
    \item \textbf{Uniform building types:} Real cities have continuous variation
    \item \textbf{Grid representation:} Square grids impose geometric artifacts
    \item \textbf{Greenfield focus:} Current model suits new cities, not retrofitting existing urban areas (requires historical preservation, renovation costs)
\end{itemize}

\subsection{Future Research Directions}

\textbf{Short-term:}
\begin{itemize}
    \item Additional building types (wind turbines, mixed-use developments)
    \item Temporal dynamics (multi-year simulation, population growth)
    \item Algorithm comparisons (reinforcement learning, simulated annealing)
\end{itemize}

\textbf{Medium-term:}
\begin{itemize}
    \item GIS data integration (import/export shapefiles, OpenStreetMap)
    \item Multi-city regional optimization
    \item Economic modeling (tax revenue, employment, housing affordability)
\end{itemize}

\textbf{Long-term:}
\begin{itemize}
    \item Digital twin integration for real-time optimization
    \item AI-assisted policy design (auto-generate zoning regulations)
    \item Extension to other SDGs (SDG 6: Clean Water, SDG 7: Clean Energy)
\end{itemize}

\subsection{Recommendations}

\textbf{For Researchers:} Prioritize curating standardized urban carbon datasets; conduct validation studies with real cities.

\textbf{For Policymakers:} Invest in monitoring systems for building-level emissions; pilot AI-driven planning in greenfield projects.

\textbf{For Practitioners:} Explore NetZero for education and preliminary analysis; provide feedback on practical constraints.

\section{Conclusion}

This paper presented NetZero, demonstrating the feasibility of Genetic Algorithms for optimizing urban layouts targeting SDG 11 and SDG 13. Through multi-objective optimization balancing carbon, happiness, and cost with eight spatial constraints, NetZero achieved near-zero carbon cities (carbon ratio $<$5\%) in 90\% of trials across diverse configurations.

\textbf{Key Contributions:}
\begin{enumerate}
    \item Comprehensive GA framework with realistic spatial constraints
    \item Empirical validation: 90\% success rate, effective trade-offs (4.1\% carbon, 78/100 happiness, \$1.05M cost)
    \item Open-source, hardware-scalable tool accessible on modest hardware
    \item Identification of critical data gaps informing future data collection efforts
\end{enumerate}

\textbf{Implications:} NetZero demonstrates that AI-driven optimization can inform sustainable city design. However, translating this to real-world impact requires: (1) standardized, region-specific emission datasets, (2) enhanced models with temporal dynamics and retrofit scenarios, (3) collaborative ecosystems involving researchers, policymakers, and planners.

While challenges remain—particularly in data availability—this research demonstrates that \textit{computational tools and methodologies exist today} to design sustainable cities. The question is no longer "Can we optimize for carbon-neutral cities?" but "Will we prioritize the data, policy, and collaboration needed to deploy these solutions at scale?"

\textbf{Code Availability:} NetZero prototype available as open-source software at [repository URL].

\section*{Acknowledgment}

This research was conducted as part of efforts to advance computational approaches to sustainable development. We acknowledge the limitations imposed by synthetic datasets and advocate for collaborative, open-data initiatives in urban carbon accounting.

\begin{thebibliography}{00}
\bibitem{UN2018} United Nations, \textit{The World's Cities in 2018—Data Booklet}, Department of Economic and Social Affairs, Population Division, 2018.

\bibitem{SDG2030} United Nations, \textit{Transforming Our World: The 2030 Agenda for Sustainable Development}, 2015. [Online]. Available: https://sdgs.un.org/2030agenda

\bibitem{Ligmann2008} A. Ligmann-Zielinska, R. L. Church, and P. Jankowski, ``Spatial optimization as a generative technique for sustainable multiobjective land-use allocation,'' \textit{Int. J. Geographical Information Science}, vol. 22, no. 6, pp. 601--622, 2008.

\bibitem{Cao2012} K. Cao, B. Huang, S. Wang, and H. Lin, ``Sustainable land use optimization using Boundary-based Fast Genetic Algorithm,'' \textit{Computers, Environment and Urban Systems}, vol. 36, no. 3, pp. 257--269, 2012.

\bibitem{Stewart2014} T. J. Stewart and R. Janssen, ``A multiobjective GIS-based land use planning algorithm,'' \textit{Computers, Environment and Urban Systems}, vol. 46, pp. 25--34, 2014.

\bibitem{Deb2002} K. Deb, A. Pratap, S. Agarwal, and T. Meyarivan, ``A fast and elitist multiobjective genetic algorithm: NSGA-II,'' \textit{IEEE Trans. Evolutionary Computation}, vol. 6, no. 2, pp. 182--197, 2002.

\bibitem{IPCC2006} IPCC, \textit{2006 IPCC Guidelines for National Greenhouse Gas Inventories}, Institute for Global Environmental Strategies (IGES), Japan, 2006.

\bibitem{EPA2022} U.S. Environmental Protection Agency, \textit{Greenhouse Gas Reporting Program (GHGRP)}, 2022. [Online]. Available: https://www.epa.gov/ghgreporting

\bibitem{Kennedy2009} C. Kennedy et al., ``Greenhouse gas emissions from global cities,'' \textit{Environmental Science \& Technology}, vol. 43, no. 7, pp. 7297--7302, 2009.

\bibitem{Seto2014} K. C. Seto et al., ``Human settlements, infrastructure and spatial planning,'' in \textit{Climate Change 2014: Mitigation of Climate Change. IPCC Working Group III Contribution to AR5}. Cambridge University Press, 2014.

\bibitem{Batty2013} M. Batty, \textit{The New Science of Cities}. MIT Press, 2013.

\bibitem{Creutzig2015} F. Creutzig, G. Baiocchi, R. Bierkandt, P. P. Pichler, and K. C. Seto, ``Global typology of urban energy use and potentials for an urbanization mitigation wedge,'' \textit{Proc. National Academy of Sciences}, vol. 112, no. 20, pp. 6283--6288, 2015.

\end{thebibliography}

\end{document}
